% answers_short.tex —— 2009 年·普通高等学校招生考试·安徽卷·理科 · 简短答案速查
% 由 tools/make_short_answers.py 自 answers.tex 的 \gksolhead/\gkanswer 提取，
% 并与 2009/2009_数学真题_安徽理_答案要点.md 逐题比对后生成，勿手改（改源 MD 后重跑）。

\begin{gknotice}
\textbf{使用说明}：本节只给最终答案便于快速核对，完整推导见解析版。答案为编写者独立推导，\textbf{非官方标准答案}。
\end{gknotice}

\gkpart{选择题}
\begin{center}\begin{tabular}{ccccccccccc}
\toprule
题号 & 1 & 2 & 3 & 4 & 5 & 6 & 7 & 8 & 9 & 10 \\
\midrule
答案 & B & D & B & A & B & C & A & C & A & D \\
\bottomrule
\end{tabular}\end{center}

\gkpart{填空题}
\begin{itemize}[leftmargin=2.6em]
  \item[\textbf{11.}] $\frac12$。
  \item[\textbf{12.}] $\sqrt{14}$。
  \item[\textbf{13.}] $127$。
  \item[\textbf{14.}] $2$。
  \item[\textbf{15.}] ①④⑤。
\end{itemize}

\gkpart{解答题}
\begin{itemize}[leftmargin=2.6em]
  \item[\textbf{16.}] （1）$\dfrac{\sqrt3}{3}$；（2）$3\sqrt2$。
  \item[\textbf{17.}] $P(X=1)=\frac13$，$P(X=2)=\frac12$，$P(X=3)=\frac16$；$EX=\frac{11}{6}$。
  \item[\textbf{18.}] （1）$90^\circ$；（2）$\dfrac{2\sqrt2}{9}$。
  \item[\textbf{19.}] $0<a\leqslant2\sqrt2$ 时在 $(0,+\infty)$ 上单调递增；$a>2\sqrt2$ 时在 $\left(0,\frac{a-\sqrt{a^2-8}}2\right)$ 与 $\left(\frac{a+\sqrt{a^2-8}}2,+\infty\right)$ 上递增，在两者之间递减。
  \item[\textbf{20.}] （1）（2）证明见解析。
  \item[\textbf{21.}] （1）证明见解析；（2）$0<a_1<1$ 或 $a_1>3$。
\end{itemize}
